% !TEX root = ../main.tex
\chapter*{凡例}\label{hanrei}\addcontentsline{toc}{chapter}{凡例}\setcounter{section}{0}\thispagestyle{plain}

\section{収録語}

本語彙集は第二著者、南琉球宮古語水納島方言の母語話者である大浦辰夫（昭和14年生）のことばを収録したものである。全部で5768項目（発音の揺れも含む）を収録した。

\section{見出し語}

見出し語はアクセント付き仮名表記、\ip{IPA}（国際音声記号）表記、品詞、意味範疇、活用、意味記述、用例、備考、発音の揺れ、類似語から構成されている。

\section{配列}

項目は仮名表記を基に五十音順に並べた（配列の際は分かち書きの空白を無視している）。従って〈ム〉から始まる語は〈む〉から始まる語と同じところにある。仮名表記が同じである場合、品詞を基準とし、接頭辞、接尾辞、助詞、終助詞、助数詞、数詞語根、連体詞、語素、名詞、動詞、形容詞、ナ形容詞、副詞、重複形、擬態語・擬音語、感動詞、接続詞、連語、諺の順で項目を並べた。

\section{仮名表記}

仮名表記は『伊良部方言辞典』の表記を踏襲した『南琉球宮古語多良間方言辞典』の表記に従った。ただし、第二著者の語感に従って有声唇歯摩擦音の二重子音\ip{[vv]}を「っヴ」に改めた。

概説で説明した通り、〈ム〉の発音が必ずしも安定していない。また、子音が続く場合、短音の〈ム〉と〈ん〉が後続の子音の調音点に同化し、完全に中和すると思われる（逆に言うとほとんどの環境において〈ム〉と〈ん〉が対立しない）。従って、本語彙集では環境で予測できない〈ム〉だけを〈ム〉と表記した。例えば、「女性」は〈みどぅム〉と発音されたため、見出し語として〈みどぅム〉を立てたが、「女子生徒」は子音が続くため、〈みどぅんしーとぅ〉のように表記した。また、項目の表記を決める際に、実際の発音によったので、同じ語根を含む複合語などの間には表記の統一を行っていない。例えば、単独の「女性」は〈みどぅム〉だが、「美人」は〈あぱらぎみどぅん〉となっている。

また、水納島方言では〈わ〉と〈ヴぁ〉の区別がないと思われるが、十分な観察ができなかったこともあり、単子音の場合は慣習的に〈わ〉を使った。

\section{アクセント}

アクセント型の認定（主に名詞、動詞、形容詞、擬態語・擬音語）は第一著者が行った。アクセント型が分かるように仮名表記にアクセント記号を付与した。平板型のa型は「\heiban{}」の記号で示した。ピッチの変動のある型は「\kako{}」と「\josho{}」の記号を用いてピッチの下降および上昇を示した。提示した下降と上昇の位置は発話のはじめに2拍の接語が続いた発音による。ただし、b型はピッチの変動が後続の接語で実現するため、語の後に「\kako{}」を付与して示した。アクセント型がピッチの変動がある発音とピッチの変動がない発音とで揺れる場合はピッチの記号を（　）で囲んだ。

各項目のアクセント型の認定は必ず複数のアクセント資料の観察に基づいて行った。名詞と擬態語・擬音語は下降調と上昇調の両方の発音を収録し、認定した。動詞は下降調における進行形と過去否定形の発音を収録し、認定した。形容詞は〈～さどぅ　あい〉、〈～さまい　ねーん〉と〈～むぬ〉の発音を収録し、認定した。十分な確信が得られなかった場合は再調査を実施してアクセント型を認定したが、アクセント型の判断が最後まで下しきれかった項目はアクセントを省略した（未調査の項目も若干数ある）。

\section{IPA表記}

音韻的な解釈に基づき、\ip{IPA}（国際音声記号）の適切な記号を選んだ。母音の無声化は音声的な特徴であると考え、表記していない。また、〈しみんやー〉\ip{[ɕiminjaː]}「セメントの家」のように、\ip{IPA}表記において音節境界が曖昧になる場合、音節境界を「\ip{.}」で示した（\ip{[ɕimin.jaː]}）。

\section{品詞}

見出し語に対し次のような品詞を設定し、括弧の中にある略号を用いて示した：助詞（助）、終助詞（終）、助数詞（助数）、数詞語根（数）、連体詞（連体）、名詞（名）、動詞（動）、形容詞（形）、ナ形容詞（ナ形）、副詞（副）、重複形（重複）、擬態語・擬音語（擬）、感動詞（感）、接続詞（接続）。それに加えて、接頭辞（接頭）、接尾辞（接尾）、語素（語素）、連語（連語）、諺（諺）の範疇も設けた。

動詞は品詞の後に活用クラス（「\rom{1}」、「\rom{2}」、「変則型」、「混合型」）も示した。活用クラスの認定は多良間方言との対応語の活用クラスを考慮しながら、過去否定形に基づいて行った。動詞によって否定形が揺れるため、活用クラスの認定を控えた項目もある。

\section{意味範疇}

一部の語に対して次の意味範疇を設けて示した：干支、貝、蟹、疑問詞（疑）、魚、指示詞（指）、植物（植）、人名、代名詞（代）、地名（地）、鳥、虫、屋号。

\section{活用}

動詞は接続形（〜して）と否定形（〜しない）の活用形を示した。活用形は見出し語に合うように、ある程度規範的に語形を整えた。

\section{意味記述}

意味は相当する共通語を示した上、必要な場合に解説を加えた。多義語は「\numbering{1}\numbering{2}\numbering{3}...」など番号を付けて意味を分けて示した。

\section{用例}

収録語の意味の理解を助けるために657点の用例を加えた。可能な限りピッチの上り・下がりを記そうとしたが、確信が得られなかった用例はピッチ記号を省略した。用例は発音通りに掲載しているため、見出し語の語形と用例における語形が一致しないことがある。

\section{備考}

必要に応じて、見出し語に関する備考を加えた。

\section{発音の揺れ}

日本語の「おととい」・「おとつい」のように同じ語に対して複数の発音が許容される場合、それぞれの発音を別の項目に立て、それぞれの項目の中に異なる発音を［同］で示し、掲載した。しかし、次の2つの場合においては異なる方針を取った。第一に、水納島方言では〈す〉〈つ〉〈ず〉が〈し〉〈ち〉〈じ〉の自由変異を持つため、原則より古い発音である前者の発音のみを立項した。例えば〈みず〉\ws{}～\ws{}〈みじ〉「水」に対して〈みず〉のみを立てた。第二に、〈さ〉行、〈ざ〉行、部分的に〈つぁ〉行が非口蓋化の発音（さ、せ、すぅ、そ）と口蓋化した発音（しゃ、しぇ、しゅ、しょ）で激しく揺れているが、一々発音の揺れを立項せず、調査の際に得られた発音をそのまま見出し語の表記として使用した。従って、本語彙集を引くに当たって、両方の文字（さ・しゃなど）で単語を探す必要がある（または共通語引きを使う）。%

\section{類似語}

見出し語と同様の意味を持つ語は［類］で示し、掲載した。

\section{日本語引き}

便宜のため、本語彙集の後に日本語引きを用意した。

\begin{table}[p]\caption*{仮名・発音記号一覧}\label{tb:kanaichiran}
\resizebox{\textwidth}{!}{    
\def\arraystretch{1}
     \begin{tabular}{|cc|cc|cc|cc|cc|cc|}
     %\hline
     \Xhline{2\arrayrulewidth}
     %\toprule
あ&\small{\ip{[a]}}&い&\small{\ip{[i]}}&& &う&\small{\ip{[u]}}&え&\small{\ip{[e]}}&お&\small{\ip{[o]}}\\
か&\small{\ip{[ka]}}&き&\small{\ip{[ki]}}&&&く&\small{\ip{[ku]}}&け&\small{\ip{[ke]}}&こ&\small{\ip{[ko]}}\\
が&\small{\ip{[ga]}}&ぎ&\small{\ip{[gi]}}&&&ぐ&\small{\ip{[gu]}}&げ&\small{\ip{[ge]}}&ご&\small{\ip{[go]}}\\
きゃ&\small{\ip{[kja]}}&&&&&きゅ&\small{\ip{[kju]}}&&&きょ&\small{\ip{[kjo]}}\\
ぎゃ&\small{\ip{[gja]}}&&&&&ぎゅ&\small{\ip{[gju]}}&&&ぎょ&\small{\ip{[gjo]}}\\
さ&\small{\ip{[sa]}}&&&す&\small{\ip{[sʉ]}}&すぅ&\small{\ip{[su]}}&せ&\small{\ip{[se]}}&そ&\small{\ip{[so]}}\\
ざ&\small{\ip{[dza]}}&&&ず&\small{\ip{[dzʉ]}}&ずぅ&\small{\ip{[dzu]}}&ぜ&\small{\ip{[dze]}}&ぞ&\small{\ip{[dzo]}}\\
しゃ&\small{\ip{[ɕa]}}&し&\small{\ip{[ɕi]}}&&&しゅ&\small{\ip{[ɕu]}}&しぇ&\small{\ip{[ɕe]}}&しょ&\small{\ip{[ɕo]}}\\
じゃ&\small{\ip{[dʑa]}}&じ&\small{\ip{[dʑi]}}&&&じゅ&\small{\ip{[dʑu]}}&じぇ&\small{\ip{[dʑe]}}&じょ&\small{\ip{[dʑo]}}\\
た&\small{\ip{[ta]}}&てぃ&\small{\ip{[ti]}}&&&とぅ&\small{\ip{[tu]}}&て&\small{\ip{[te]}}&と&\small{\ip{[to]}}\\
だ&\small{\ip{[da]}}&でぃ&\small{\ip{[di]}}&&&どぅ&\small{\ip{[du]}}&で&\small{\ip{[de]}}&ど&\small{\ip{[do]}}\\
ちゃ&\small{\ip{[tɕa]}}&ち&\small{\ip{[tɕi]}}&つ&\small{\ip{[tsʉ]}}&ちゅ&\small{\ip{[tɕu]}}&ちぇ&\small{\ip{[tɕe]}}&ちょ&\small{\ip{[tɕo]}}\\
つぁ&\small{\ip{[tsa]}}&&&&&つぅ&\small{\ip{[tsu]}}&つぇ&\small{\ip{[tse]}}&つぉ&\small{\ip{[tso]}}\\
な&\small{\ip{[na]}}&に&\small{\ip{[ni]}}&&&ぬ&\small{\ip{[nu]}}&ね&\small{\ip{[ne]}}&の&\small{\ip{[no]}}\\
にゃ&\small{\ip{[nja]}}&&&&&にゅ&\small{\ip{[nju]}}&&&にょ&\small{\ip{[njo]}}\\
は&\small{\ip{[ha]}}&ひ&\small{\ip{[hi]}}&&&ほぅ&\small{\ip{[hu]}}&へ&\small{\ip{[he]}}&ほ&\small{\ip{[ho]}}\\
ば&\small{\ip{[ba]}}&び&\small{\ip{[bi]}}&&&ぶ&\small{\ip{[bu]}}&べ&\small{\ip{[be]}}&ぼ&\small{\ip{[bo]}}\\
ぱ&\small{\ip{[pa]}}&ぴ&\small{\ip{[pi]}}&&&ぷ&\small{\ip{[pu]}}&ぺ&\small{\ip{[pe]}}&ぽ&\small{\ip{[po]}}\\
ひゃ&\small{\ip{[hja]}}&&&&&ひゅ&\small{\ip{[hju]}}&&&ひょ&\small{\ip{[hjo]}}\\
びゃ&\small{\ip{[bja]}}&&&&&びゅ&\small{\ip{[bju]}}&&&びょ&\small{\ip{[bjo]}}\\
ぴゃ&\small{\ip{[pja]}}&&&&&ぴゅ&\small{\ip{[pju]}}&&&ぴょ&\small{\ip{[pjo]}}\\
ふぁ&\small{\ip{[fa]}}&ふぃ&\small{\ip{[fi]}}&&&ふ&\small{\ip{[fu]}}&ふぇ&\small{\ip{[fe]}}&ふぉ&\small{\ip{[fo]}}\\
ヴぁ&\small{\ip{[va]}}&ヴぃ&\small{\ip{[vi]}}&&&ヴぅ&\small{\ip{[vu]}}&ヴぇ&\small{\ip{[ve]}}&ヴぉ&\small{\ip{[vo]}}\\
ま&\small{\ip{[ma]}}&み&\small{\ip{[mi]}}&&&む&\small{\ip{[mu]}}&め&\small{\ip{[me]}}&も&\small{\ip{[mo]}}\\
みゃ&\small{\ip{[mja]}}&&&&&みゅ&\small{\ip{[mju]}}&&&みょ&\small{\ip{[mjo]}}\\
や&\small{\ip{[ja]}}&&&&&ゆ&\small{\ip{[ju]}}&いぇ&\small{\ip{[je]}}&よ&\small{\ip{[jo]}}\\
ら&\small{\ip{[ra]}}&り&\small{\ip{[ri]}}&&&る&\small{\ip{[ru]}}&れ&\small{\ip{[re]}}&ろ&\small{\ip{[ro]}}\\
っら&\small{\ip{[ɭɭa]}}&&&&&&&&&っろ&\small{\ip{[ɭɭo]}}\\
りゃ&\small{\ip{[rja]}}&&&&&りゅ&\small{\ip{[rju]}}&&&りょ&\small{\ip{[rjo]}}\\
わ&\small{\ip{[ʋa]}}&&&&&&&&&を&\small{\ip{[ʋo]}}\\
ん&\small{\ip{[n]}}&ム&\small{\ip{[m]}}&ヴ&\small{\ip{[v]}}&&&&&&\\
ー&\small{\ip{[ː]}}&\multicolumn{3}{c}{っ\hspace{2.5pt}（子音を重ねる）}&&&&&&&\\
     %\hline
     %
     \Xhline{2\arrayrulewidth}
     %\bottomrule
 \end{tabular}
 }
\end{table}

